\documentclass[11pt]{article}
\usepackage{fontspec}
\setmainfont{Times New Roman}
\setsansfont{Arial}
\setmonofont{Consolas}
\usepackage{amsmath,amssymb,mathtools}
\usepackage{geometry}
\usepackage{microtype}
\geometry{a4paper,margin=2.05cm}
\setlength{\parindent}{1.1em}
\setlength{\parskip}{0pt}
\providecommand{\frakO}{\mathfrak{O}}
\providecommand{\frako}{\mathfrak{o}}
\providecommand{\frakp}{\mathfrak{p}}
\providecommand{\Gg}{\mathfrak{G}}
\emergencystretch=220em
\allowdisplaybreaks
\sloppy
\hbadness=10000
\vbadness=10000
\hfuzz=4pt
\vfuzz=10pt
\raggedbottom
\begin{document}% Unit: paper37_source_fidelity_v001.
% Language lane: Interslavic.
% Direct source: sources/paper37/source_fidelity/Noether_Paper37_ORIGINAL_JReineAngewMath167_R124plus_source_fidelity_witness_v001.tex.
\section*{37. Normalna baza v poljah bez vyšego razvětvjenja}

\begin{center}
\emph{Journal f. d. reine u. angew. Math. 167 (1932), s. 147--152}
\end{center}

Pod normalnoju bazoju Galoisovogo čislovogo polja \(K/k\) razuměje se baza maksimalnogo porjadka \(\frakO\) polja \(K\) vzhodno k maksimalnomu porjadku \(\frako\) polja \(k\), sostojęča iz konjugovanyh elementov. Nužno uslovje za bytje takoj bazy jest poznano\footnote{Srav. A. Speiser, Gruppendeterminante und Körperdiskriminante, § 6. Math. Ann. 77 (1916), s. 546--562.}: v \(K/k\) ne smųt nastupati vyša polja razvětvjenja. To ostaje pravilno i kogda se ograničimo na jedno město, to jest kogda prěhodimo k \(\frakp\)-adičnomu razširjenju \(k_\frakp\) polja \(k\) i k odgovornomu \(K_\frakp\) polja \(K\). Niže pokazujem obratno tvrdženje:

\emph{Za vsa města \(\frakp\), kde \(\frakp\) ne děli stepen \(K/k\), jest normalna baza. Osoblivo: ako \(K/k\) ne ima vyšego razvětvjenja, togda pri vsakom razvětvjenom městu jest normalna baza; diskriminant tam stava kvadratom grupovej determinanty.}

K tomu poslědnjemu tvrdženju treba dodati, že, protivno predpostavjenju E. Artina, prěhod k jego provodnikam\footnote{E. Artin, Über die gruppentheoretische Struktur der Diskriminante. J. f. Math. 164 (1931), s. 1--11.} potrěbuje ještě daljne razmysljenja: razklad grupovej determinanty po ireducibilnyh karakterah daje obobčene koreneve čisla; odgovorno sobranje daje razklad v osnovnom polju, razširjenom karakterami. V najprostějših slučajah mogu te nove faktory identificirati s Artinovymi provodnikami, pomoću idealno-teoretičnogo razklada korenevyh čisel.\footnote{[6. 10. 31.] Že i občno budų potrebne ještě idealno-teoretične razmysljenja, pokazuje slědujuči, mi od M. Deuringa soobčeny i voljno varijovany, primjer nemaksimalnogo porjadka, jegov diskriminant možno predstaviti kako kvadrat grupovej determinanty, dok konjugovanym karakteram odgovarjajų različne idealne faktory. Nehaj \(k\) jest polje petyh korenjev jediničnosti, \(K=k(\sqrt[5]{2})\). Razgleda se porjadok \([1,2\sqrt[5]{2},\sqrt[5]{2^2},\sqrt[5]{2^3},\sqrt[5]{2^4}]\), imenno pri městu \((2)\), kde on po razmysljenjah toj noty ima normalnu bazu. Razklad po karakterah daje za odgovornu grupovu determinantu upravo vyšej navedene bazove elementy kako faktory. Zato ``provodniki'', izključivši \(1\), stavajų: \(2\sqrt[5]{2}/\sqrt[5]{2^4}\), \(\sqrt[5]{2^2}/\sqrt[5]{2^3}\), \(\sqrt[5]{2^3}/\sqrt[5]{2^2}\), \(\sqrt[5]{2^4}/(2\sqrt[5]{2})\), t. j. \(4,2,2,4\), i tako sut različne za konjugovane karaktere.} Dalje hoču primětiti, že i v občnoj situaciji, kde normalna baza ne obstaje, razčlenjenje na obobčene koreneve čisla jest vsegda možno pomoću ``kompozicijnogo reda po Galoisovym modulam''; i že iz togo analogno kako gore vyhodi razklad v razširjenom osnovnom polju; tu znovu počinajut idealno-teoretična vprašanja.

Dokaz byta normalnej bazy pod gornymi uslovjami osnovan jest na tom, že \(\frakO/\frako\), gledano kako Galoisov modul, to jest kako \(\frako\)-modul s podstanovami Galoisovej grupy kako oblastju operatorov, stava operatorno izomorfny k jednostrannomu idealu celočislnogo grupovogo kolca, razširjenogo s \(\frako\). To kolco v vsih običnyh razvětvjenih městah jest maksimalny porjadok; ale idealy v maksimalnyh porjadkah \(\frakp\)-adičnyh poluprostyh sistemov sut glavne idealy\footnote{Srav. H. Hasse, Über \(\frakp\)-adische Schiefkörper und ihre Bedeutung für die Arithmetik hyperkomplexer Zahlsysteme, teoremy 46 i 47. Math. Ann. 104 (1931), s. 495--534. Pri komutativnyh sistemah svojstvo glavnyh idealov plati uže pri prěhodu k kvocientnomu kolcu po \(\frakp\) v \(k\); zato možno se ograničiti na to slabše razširjenje za definiciju města, kogda ide o abelove grupy i polja.}, zato tu nastavajut iz jednogo elementa podstanovami grupy ili grupovogo kolca. Vračajuči se ot togo k Galoisovomu modulu, dostajemo ravno bytje normalnej bazy. Pri městah vyšego razvětvjenja celočislno grupovo kolco prěhodi v nemaksimalny porjadok, a ideal, odgovorny \(\frakO/\frako\), v neglavny ideal.

Vsa ta razmysljenja očividno ostavajut pravilne, ako ide o funkcijska polja jednoj proměnnoj, tako že karakteristika konstantnogo polja ne děli stepen \(K/k\).

\subsection*{§1. \(\frakp\)-adično razširjeno celočislno grupovo kolco}

Nehaj \(\frakO\), respektivno \(\frako_\frakp\), označajut maksimalne porjadky algebraičnogo čislovogo polja \(k\) i jego \(\frakp\)-adičnogo razširjenja \(k_\frakp\), kde \(\frakp\) jest prosty ideal iz \(k\). Dalje nehaj \((\Gg)\) znači racionalno, a \([\Gg]\) celočislno grupovo kolco grupy \(\Gg\) s čislom elementov \(n\). Pod \((\Gg)_k\), respektivno \((\Gg)_{k_\frakp}\), razuměje se grupovo kolco s koeficientami iz \(k\), respektivno iz \(k_\frakp\); podobno pod \([\Gg]_\frako\), respektivno \([\Gg]_{\frako_\frakp}\), razuměje se razširjenje celočislnogo grupovogo kolca k takomu s koeficientami iz \(\frako\), respektivno iz \(\frako_\frakp\).

Tak \((\Gg)_k\) i \((\Gg)_{k_\frakp}\) sut sistemy bez radikala, to jest poluproste sistemy, a \([\Gg]_\frako\) i \([\Gg]_{\frako_\frakp}\) sut porjadky iz \((\Gg)_k\), respektivno iz \((\Gg)_{k_\frakp}\).

\paragraph{Teorema 1.}
\emph{Ako prosty ideal \(\frakp\) ne děli čislo elementov \(n\) grupy \(\Gg\), togda \([\Gg]_{\frako_\frakp}\) stava maksimalny porjadok v \((\Gg)_{k_\frakp}\); ako \(\frakp\) děli \(n\), togda \([\Gg]_{\frako_\frakp}\) ne jest maksimalny.}

Diskriminant celočislnogo grupovogo kolca, tak že i diskriminant \([\Gg]_\frako\) i \([\Gg]_{\frako_\frakp}\), jest ravny někoj stepeni čisla elementov \(n\).\footnote{Srav. napr. E. Noether, Hyperkomplexe Größen und Darstellungstheorie, § 26. Math. Ztschr. 30 (1929), s. 641--692.} Zato ako \(\frakp\) ne děli \(n\), diskriminant \([\Gg]_{\frako_\frakp}\) jest jedinica. To znači, že \([\Gg]_{\frako_\frakp}\), pri fiksovanom \(\frako_\frakp\), ne možno razširiti do obširnějšego porjadka; takomu razširjenju odgovarjalo by odděljenje kvadratnogo nej ediničnogo faktora od diskriminanta. Ale \(\frako_\frakp\) uže bylo predpostavljeno kako maksimalny porjadok, to jest maksimalny v \(k_\frakp\). Tak dokazana jest prva čest teoremy 1, jedina upotrebjena niže.

Ako protivno \(\frakp\) děli \(n\), togda direktna suma, odgovorna jediničnomu predstavjenju,
\[
        \mathfrak C=E^{(1)}[\Gg]_{\frako_\frakp}+(1-E^{(1)})[\Gg]_{\frako_\frakp},
        \qquad
        E^{(1)}=\frac1n\sum_{S\in\Gg} S,
\]
jest istinno razširjenje \([\Gg]_{\frako_\frakp}\). Bo iz \(E^{(1)}=\frac1n\sum S\) slěduje, že \(\mathfrak C\) jest istinno razširjenje; \(\mathfrak C\) jest porjadok so zavisnymi poroditeljami \(E^{(1)}\), \((1-E^{(1)})S\) za \(S\in\Gg\), pri čem \(E^{(1)}S=E^{(1)}\).

\paragraph{Teorema 2.}
\emph{Ako \(\frakp\) ne děli \(n\), togda vsak cěly ili drobny ideal vzhodno k \([\Gg]_{\frako_\frakp}\) jest glavny ideal.}

Bo po teoremi 1 \([\Gg]_{\frako_\frakp}\) jest maksimalny porjadok \(\frakp\)-adičnogo poluprostogo sistema; zato idealy sut glavne idealy.\footnote{Hasse, navedeno dělo. Svojstvo glavnyh idealov tam dokazano jest samo za proste sistemy, ale iz togo po poznanyh zaključkah slěduje ono i za poluproste. Komponenty maksimalnogo porjadka sut znovu maksimalne; komponenty gledanogo ideala zato sut glavne idealy, recimo s bazami \(a_i\). Togda \(a=\sum a_i\) jest baza izhodnogo ideala. Za abelove grupy srav. još primětku 3.}

\subsection*{§2. Galoisove moduly, operatorna izomorfija, normalna baza}

Nehaj \(K/k\) jest Galoisovo, a \(\Gg\) grupa; \(z^S\) neka znači element \(K\), nastaly iz \(z\) podstanovoju \(S\) iz \(\Gg\). Podstanovy \(\Gg\) porađajut na \(K/k\), gledanom kako \(k\)-modul, to jest aditivna abelova grupa s \(k\) kako oblastju operatorov, operatorne avtomorfizmy. Zato možno \(K/k\) definirati kako modul vzhodno k \((\Gg)_k\) pravilo:
\[
        z\Bigl(\sum_i S_i c_i\Bigr)=\sum_i z^{S_i}c_i,
        \qquad z\in K,\; S_i\in\Gg,\; c_i\in k.
\]

\paragraph{Definicija.}
\emph{Vsak \((\Gg)_k\)-modul iz \(K/k\) imenuje se racionalny Galoisov modul. Podobno \([\Gg]_\frako\)-moduly iz \(K/k\) imenujut se cěle Galoisove moduly; osoblivo \(\frakO/\frako\) jest cěly Galoisov modul.}

\paragraph{Teorema 3.}
\emph{Kako racionalny Galoisov modul, \(K/k\) jest operatorno izomorfny k \((\Gg)_k\).}

To slěduje iz togo, že \(K/k\) vsegda ima poljnu normalnu bazu, to jest \(k\)-bazu iz konjugovanyh elementov \(z^S\).\footnote{Bo ako \(a_1,\ldots,a_n\) jest baza \(K/k\) i \(u_i\) označajut neizvěstne, togda \(\sum u_i a_i^S\) sut linejno nezavisne, kogda \(S\) prohodi elementy \(\Gg\); zato možno \(u_i\) specializovati k elementam iz \(k\) tako, že linejna nezavisnost ostane.} Priradenje
\[
        S\longmapsto z^S,
        \qquad
        \sum_i S_i c_i\longmapsto \sum_i z^{S_i}c_i
        \quad(c_i\in k),
        \quad\hbox{to jest } ST\longmapsto z^{ST},
\]
daje operatorny homomorfizm, ktory zaradi ravnogo ranga nad \(k\) stava izomorfizm.

\paragraph{Teorema 4.}
\emph{Kako cěly Galoisov modul, \(\frakO/\frako\) jest operatorno izomorfny k jednomu \([\Gg]_\frako\)-modulu iz \((\Gg)_k\), maksimalnogo ranga \(n\). Podobno \(\frakO_\frakp/\frako_\frakp\) jest operatorno izomorfny k \([\Gg]_{\frako_\frakp}\)-modulu ranga \(n\) iz \((\Gg)_{k_\frakp}\).}

\((\Gg)_k\)-izomorfizm \(K/k\to(\Gg)_k\), dany v teoremi 3, priradžuje \([\Gg]_\frako\)-modulu \(\frakO/\frako\) \([\Gg]_\frako\)-modul ranga \(n\) iz \((\Gg)_k\). Dalje možno toj \((\Gg)_k\)-izomorfizm razširiti do \((\Gg)_{k_\frakp}\)-izomorfizma od \(K_\frakp/k_\frakp\) k \((\Gg)_{k_\frakp}\), ako samo pustimo koeficienty \(c_i\) v priradenju teoremy 3 proběgati vse elementy iz \(k_\frakp\). Toj izomorfizm potom induciruje \([\Gg]_{\frako_\frakp}\)-izomorfizm \(\frakO_\frakp/\frako_\frakp\). Pri tom treba \(K_\frakp/k_\frakp\) gledati kako hiperkompleksny sistem nad \(k_\frakp\). On nastava iz \(K/k\) razširjenjem koeficientov; občno stava sistem s nulovymi děliteljami, imenno suma izomorfnyh polj, zato poluprosty sistem.

\paragraph{Teorema 5.}
\emph{Pri vsakom městu \(\frakp\), ktore ne děli stepen \(n\) od \(K/k\), \(\frakO_\frakp/\frako_\frakp\) ima normalnu bazu.}\footnote{V slučaju abelovyh polj možno pri tom město po primětkah 3 i 5 definirati takože kvocientnym kolcom.}

Bo po teoremi 1 \([\Gg]_{\frako_\frakp}\) tam stava maksimalny porjadok; \([\Gg]_{\frako_\frakp}\)-moduly ranga \(n\) iz \((\Gg)_{k_\frakp}\) zato stavajut cěle ili drobne idealy, imenno glavne idealy po teoremi 2. Ako tedy ideal, odgovorny \(\frakO_\frakp/\frako_\frakp\), jest \(W[\Gg]_{\frako_\frakp}\), togda on ima \(\frako_\frakp\)-bazu \(WS_i\); operatorna izomorfija govori, že \(w^{S_i}\) jest \(\frako_\frakp\)-baza \(\frakO_\frakp/\frako_\frakp\), ako \(w\) označa element iz \(\frakO_\frakp/\frako_\frakp\), odgovorny \(W\). Tih \(n\) konjugovanyh elementov \(w^{S_i}\) zato tvori normalnu bazu \(\frakO_\frakp/\frako_\frakp\).

Ako \(\frakp\) děli \(n\), odgovorny modul \(\frakO_\frakp/\frako_\frakp\) ne može biti glavny modul, bo v tom slučaju normalna baza ne može obstajati.

\paragraph{Dodatna primětka k teoremi 3.}
Iz operatornej izomorfije \((\Gg)_k\) k \(K/k\), dokazanej v teoremi 3,\footnote{Dokaz očividno predpostavja samo, že \(K\) jest Galoisovo prvogo roda nad \(k\), i že \(k\) ima beskonečno mnogo elementov. To poslědnje ograničenje jest nepotrebno: M. Deuring našel jest dokaz teoremy 3, valjany takože za polja s konečno mnogo elementov; iz jego obratno slěduje bytje poljnej normalnej bazy. Jednovremenno iz togo vyhodi dokaz byta primitivnogo elementa, ktory jednako rabota za polja s konečno i beskonečno mnogo elementov.} slědujut Speiserove rezultaty o Kleinovom problemu form,\footnote{Navedeno dělo. Pojem Galoisovogo modula abstrahovan jest odtud. Ako karakteristika \(k\) děli \(n\), togda ide o kompozicijne redy po Galoisovym modulam i jih ireducibilne faktory.} ktore možno formulovati tako:

\emph{Vsako nad \(k\) ireducibilno predstavjenje Galoisovej grupy \(\Gg\) od \(K/k\) porađano jest ireducibilnym racionalnym Galoisovym modulom iz \(K/k\); vsako absolutno ireducibilno predstavjenje porađano jest Galoisovym modulom iz \(K_Z/Z\).}

Tu \(Z\) znači razkladno polje, obnimajuče \(k\), za odgovornu komponentu grupovogo kolca; \(K_Z/Z\) treba gledati kako hiperkompleksny sistem nad \(Z\), nastaly iz \(K/k\) razširjenjem koeficientov. Osoblivo \(K_Z\) ostaje polje, ako možno \(Z\) izbrati prosto k \(K\) vzhodno k \(k\), to jest ako \(K_Z/Z\) stava akcesorno razširjenje, slučaj obrabotany Speiserom.

Samo tvrdženje slěduje pomoću operatornej izomorfije iz odgovornogo tvrdženja za \((\Gg)_k\), respektivno \((\Gg)_Z\), kde ireducibilne idealy porađajut predstavjenje.

Ako \(v_1,\ldots,v_l\) jest baza takogo Galoisovogo modula, a \(S\mapsto \bar S\) odgovorno predstavjenje, togda:
\[
        \begin{pmatrix}\vdots\\ v_i^S\\ \vdots\end{pmatrix}
        =\bar S\begin{pmatrix}\vdots\\ v_i\\ \vdots\end{pmatrix}.
\]
To jest formulacija Kleina i Speisera.

\subsection*{§3. Diskriminant kako grupova determinanta v poljah bez vyšego razvětvjenja}

Za razklad diskriminanta dostatočno jest, kako poznano, gledati razklad pri jednotlivyh razvětvjenih městah; v poljah bez vyšego razvětvjenja ide zato samo o města s normalnoju bazoju.

\paragraph{Teorema 6.}
\emph{Pri Galoisovyh poljah \(K/k\) bez vyšego razvětvjenja diskriminant pri vsakom razvětvjenom městu jest ravny kvadratu grupovej determinanty Galoisovej grupy, s neizvěstnymi zaměnjenymi normalnoju bazoju. Odgovorno ireducibilnym predstavjenjam grupova determinanta razpada se na obobčene koreneve čisla; diskriminant razpada se na idealy osnovnogo polja, razširjenogo karakterami, pri čem konjugovano-kompleksnym karakteram odgovarjajut ty same idealy.}

Bo s \(w^S\) kako normalnoju bazoju diskriminant \(\Delta\) jest kvadrat od
\[
        D=\bigl|w^{ST^{-1}}\bigr|,\qquad S,T\in\Gg.
\]
\(D\) jest grupova determinanta s \(w^S\) na městu neizvěstnyh \(x_S\).

Razklad v koreneve čisla vyhodi iz zaključkov Speisera, navedeno dělo. Nehaj
\[
        D=D_1^{f_1}\cdots D_t^{f_t},
        \qquad
        M=f_1M_1+\cdots+f_tM_t
\]
jest razklad grupovej determinanty, respektivno grupovej matrice, po absolutno ireducibilnyh predstavjenjah. Ako \(S\mapsto\bar S\) znači predstavjenje, odgovorno \(M_\lambda\), gde \(\bar S\) leže v razširjenom osnovnom polju, togda dostajemo:
\[
        M_\lambda=\sum_S w^S\bar S,
        \qquad
        M_\lambda^{T^{-1}}=\sum_S w^{ST^{-1}}\bar S
          =\sum_R w^R\bar R\bar T=M_\lambda\bar T,
\]
i prěhodom k determinanti
\[
        D_\lambda^{T^{-1}}=D_\lambda|\bar T|=D_\lambda\varepsilon_T.
\]
Tako \(D_\lambda\) poznavajut se kako obobčene koreneve čisla; \(\varepsilon_T\), kako determinanta \(\bar T\), jest korenj jediničnosti, to jest ireducibilno predstavjenje faktor-grupy \(\Gg\) po komutatornoj grupě. V slučaju cikličnyh grup \(D_\lambda\) sut prosto Lagrangeova koreneva čisla \(\sum w^S\chi_\lambda(S)\).

Občno \(D_\lambda\) leže v hiperkompleksnom sistemu, nastalom iz \(K_\frakp/k_\frakp\) tako, že oblast koeficientov \(k_\frakp\) razširjena jest odgovornym karakterom; i imenno ide o \emph{cěle} veličiny togo sistema. Bo determinanta, sostavjena s neizvěstnymi \(x_S\), ima koeficienty, ktore sut cěle čisla iz polja karaktera; a \(w^S\) sut cěle elementy iz \(K_\frakp\).

Da by se prěšlo od razklada grupovej determinanty k razkladu diskriminanta, vsaky raz se predstavjenje sobiraje s jego adjungovanym. Ako \(\lambda,\bar\lambda\) označajut adjungovane predstavjenja, togda jednočasno:
\[
        D_\lambda^{T^{-1}}=D_\lambda\varepsilon_T,
        \qquad
        D_{\bar\lambda}^{T^{-1}}=D_{\bar\lambda}\varepsilon_T^{-1};
\]
jednovremenno determinanty, sostavjene za neizvěstne \(x_S\) i prinadležne k \(\lambda,\bar\lambda\), sut konjugovano-kompleksne; medžu tym občno konjugovanym karakteram pri neizvěstnyh \(x_S\) odgovarjajut konjugovane determinanty.

Položimo zato \(\Delta_\lambda=D_\lambda D_{\bar\lambda}\). Togda treba adjungovati samo realno podpolje \(\lambda\)-togo karaktera, i plati
\[
        \Delta_\lambda^T=\Delta_\lambda
        \qquad\hbox{za vse }T\hbox{ iz }\Gg;
\]
i zato\footnote{Poneže ide o hiperkompleksno razširjenje koeficientov polja, obyčny zaključok Galoisovej teorije treba modifikovati. Nehaj \(P\) jest gledano razširjeno polje od \(k_\frakp\) i
\[
        (K_\frakp)_P=(K_\frakp)_\mathfrak P e^{S_1}+\cdots+(K_\frakp)_\mathfrak P e^{S_r}
\]
jest direktny razklad v polja. Togda \((K_\frakp)_\mathfrak P e^{S_i}/Pe^{S_i}\) jest Galoisovo, grupa jest podgrupa razkladnoj grupy jednogo prostogo faktora od \(\frakp\), i \(e^{S_i}\) sut konjugovane vzhodno k pobočnym klasam. Iz \(\Delta_\lambda^T=\Delta_\lambda\) najprvo slěduje, že \(i\)-ta komponenta \(\Delta_\lambda\) leži v \(Pe^{S_i}\), zato jest recimo \(\gamma_i e^{S_i}\) s \(\gamma_i\) v \(P\). Dalje iz konjugovanosti \(e^{S_i}\) slěduje, že vsi \(\gamma_i\) sut ravni, zato vpravdě \(\Delta_\lambda=\gamma\sum e^{S_i}=\gamma\) leži v \(P\).}: \(\Delta_\lambda\) leži v osnovnom polju, razširjenom realnym podpoljem \(\lambda\)-togo karaktera.

\emph{Razklad diskriminanta}
\[
        \Delta=\Delta_1^{f_1}\cdots\Delta_t^{f_t}
\]
\emph{jest zato takov razklad v osnovnom polju, razširjenom realnymi podpoljami karakterov, pri čem konjugovano-kompleksnym karakteram odgovarjajut ty same faktory. Za ostale konjugovane karaktere bez daljnego nič ne slěduje.}

Tym teorema 6 dokazana jest vo vsih čestjah. Ako možno pokazati, že idealy, porađene \(\Delta_\lambda\), vpravdě uže leže v osnovnom polju \(k\) i za konjugovane karaktere stavajut ravne, togda one sovpadajut s Artinovymi provodnikami, kogda samo složenym karakteram priradimo odgovorne stepenne produkty \(\Delta_\lambda\). Bo plati i funkcionalno uravnenje i fakt, ktory direktno slěduje iz normalnej bazy (srav. Speiser, navedeno dělo): jediničnym predstavjenjam podgrup odgovarjajut diskriminanty odgovornyh podpolj. Iz sovpadanja pri vsakom městu slěduje sovpadanje polnyh idealov. Ako se ograničimo na racionalno ireducibilne karaktere, navedene fakty dajut sovpadanje. Za absolutno ireducibilne karaktere bude sovpadanje v najprostějšem slučaju još direktno potvrdženo:

\paragraph{Teorema 7.}
\emph{Ako \(k\) jest polje racionalnyh čisel, \(K/k\) ciklično prostogo stepena \(l\), prostogo k diskriminantu, zato bez vyšego razvětvjenja, togda idealy, porađene \(\Delta_\lambda\) za nejedinične predstavjenja, sut vsi ravne medžu soboju i sovpadajut s provodnikom \(K/k\).}

To slěduje iz poznanogo razklada korenevyh čisel.\footnote{Hilbert, Zahlbericht, § 108. Poneže po teoremi 5 bytje normalnej bazy jest ustanovjeno, a město po primětki 7 možno tu prosto definirati kvocientnym kolcom, ne treba više upotrěbiti fakt iz Zahlberichta, že absolutno-ciklična polja sut kružna polja.} Pri razvětvjenom městu \(\frakp\) od \(K/k\) plati za \(k_\varepsilon\), polje \(l\)-tyh korenjev jediničnosti, i za \(K_\varepsilon\) (\(K_\varepsilon\) ostaje polje, bo \(k_\varepsilon\) jest prosto k \(K\) i \(\frakp\)-adično razširjenje po primětki 7 tu jest nepotrebno):
\[
        (p)=\frakp_1\cdots\frakp_{l-1},
        \qquad
        \frakp_i=\mathfrak P_i^{\,l},
\]
s \(\frakp_i\) v \(k_\varepsilon\), \(\mathfrak P_i\) v \(K_\varepsilon\). Dalje za koreneve čisla plati:
\[
        (\Omega_\lambda)=\mathfrak P_1^{r_1}\cdots\mathfrak P_{l-1}^{r_{l-1}},
        \qquad
        (\Omega_{\bar\lambda})=
        \mathfrak P_1^{l-r_1}\cdots \mathfrak P_{l-1}^{l-r_{l-1}},
\]
kde \(r_1,\ldots,r_{l-1}\) jest permutacija čisel \(1,\ldots,l-1\). Koreneve čisla \(\Omega_\lambda\) sut ale tu identične s faktorami \(D_\lambda\) grupovej determinanty; zato:
\[
        (\Delta_\lambda)=(D_\lambda D_{\bar\lambda})
        =(\Omega_\lambda\Omega_{\bar\lambda})
        =\mathfrak P_1^{l}\cdots\mathfrak P_{l-1}^{l}
        =\frakp_1\cdots\frakp_{l-1}=(p);
\]
zato jest sovpadanje s provodnikom \(K/k\).

\begin{center}
Prijeto 24 avgusta 1931.
\end{center}
\end{document}